\section{Completed recursive families and a total envelope}
\label{jfam:section}

The geometric doubling mechanism belongs to BBL~\cite{bbl2008}. A compatible seed supplies a tangent grid, simplicity, consecutive axis caps, positive central apex and certified triangle list. A visible seed additionally supplies an admissible normal and rightmost boundary pair.

\begin{theorem}[Verified BBL step]\label{jfam:doubling}
Let $q=4r\ge20$, $0<\varepsilon<\tan(\pi/(2q))$, and $t_j=\tan(j\pi/q)$. Take $Y_0:y=0$ and $L_j:y=m_j(x-a_j)$ on
\[
 (a_j)=(-t_{q/2-1},\ldots,-t_1,-\varepsilon,\varepsilon,t_1,\ldots,t_{q/2-1}).
\]
If these lines are simple, all $q-1$ axis caps are triangular, the central apex is positive and $T$ distinct triangles are certified, then $\Ks(2q+1)\ge T+q^2$.
\end{theorem}
\begin{proof}[Geometric mechanism]
For $\beta_i=-\pi/2+(i+1/2)\pi/q$, add
$y=\kappa(\sin(2\beta_i)+\delta\cot\beta_i)(x-\tan\beta_i)$.
Choose $\delta>0$, then $\kappa>0$, sufficiently small. Exact crossing order yields $2\binom{q+1}{2}-q=q^2$ mixed triangles. Old triangles off $Y_0$ persist, alternating caps replace its noncentral triangles, and the central triangle survives. The step returns the same geometric invariants, after actual support relabelling. A visible seed also retains projection order and boundary data.
\end{proof}

The eleven-line realization has reciprocal slopes
$(15,-2,14,1,4,-4,5,-3,12,-5)/10$ and, on $0<\varepsilon\le10^{-5}$, 32 triangles, nine caps and five visible pairs. Its count comes from a Honma-based input~\cite{savchuk2025}. A separate 21-line uniform seed has 132 triangles, 19 caps and ten visible pairs and starts the $r\ge5$ induction.

\begin{theorem}[Retained infinite families]\label{jfam:recursive}
For $t\ge0$, put $q=10\cdot2^t$ and $T_q=(q^2-4)/3$. Then
\begin{equation}\label{jfam:families}
 \Ks(q+1)\ge T_q,\qquad \Ks(q+2)\ge T_q+q/2.
\end{equation}
\end{theorem}
\begin{proof}
Uniform availability on some positive parameter interval survives every step. Starting with $(r,T,V)=(5,132,10)$, the actual transformation
$(r,T,V)\mapsto(2r,T+(4r)^2,V+2r)$
gives the displayed counts and $V=q/2$. The exterior theorem supplies the even companion. The eleven-line seed covers $t=0$. The interval may shrink with depth; a common positive interval for all depths is not claimed.
\end{proof}

\begin{theorem}[One-triangle simple windows]\label{jfam:windows}
For the same $q$,
\[
 T_q\le\Ks(q+1)\le T_q+1,\qquad
 T_q+q/2\le\Ks(q+2)\le T_q+q/2+1.
\]
\end{theorem}
\begin{proof}
Apply the actual simple bounds of Section~\ref{jup:section}; their floors are $T_q+1$ and $T_q+q/2+1$.
\end{proof}
The upper proofs use standard axioms. Existence inherits four finite native checks for the odd source and six for its even companion. This is a simple-model window; multiple intersections remain allowed in $\K$.

\begin{theorem}[Compatible 61-line orbit]\label{jfam:61}
For every $t\ge0$, put $q=60\cdot2^t$. Actual simple arrangements give
\begin{equation}\label{jfam:61-formula}
 \Ks(q+1)\ge1200\,4^t-10,\qquad
 \Ks(q+2)\ge1200\,4^t+30\,2^t-11.
\end{equation}
Their gaps to the respective simple upper floors are nine and ten.
\end{theorem}
\begin{proof}
A reciprocal-slope refit of Rohith Poola's OpenMath input~\cite{poola2026} preserves 1190 triangles and 59 axis caps for every real $0<\varepsilon\le10^{-8}$ on the actual tangent60 grid. Exact tangent enclosures and finite sign checks establish simplicity and all seed hypotheses. A further cap-cone mutation supplies 29 visible pairs, including the required boundary pair. BBL iteration gives $V_t=30\cdot2^t-1$, and exterior realization gives the even formula. See
\proofref{Kobon/OpenMathConstructionParetoFamily61.lean}{48}{\lean{odd_family}} and
\proofref{Kobon/OpenMathConstructionParetoFamily61.lean}{51}{\lean{even_family}}.
\end{proof}
Poola supplied the finite count 61:1190; BBL supplied doubling. The uniform compatible refit and 29-pair improvement are the present developments, with seven explicit finite native roots in the even proof. The earlier 28-pair seed remains preserved; its even count is one smaller at every depth. At 62 lines the new 1219 is below the retained $\G(62)=1220$. For $t\ge1$, the odd/even gains over $\G$ are $20\cdot2^t-11$ and $10\cdot2^t-11$: for example $121:4790$, $122:4849$, $241:19190$, $242:19309$. These are repository comparisons, without absolute numerical-priority claims.
\begin{figure}[tb]
\centering\includegraphics[width=.96\linewidth]{new-family61-gaps.pdf}
\caption{Exact formulas for $q=60\cdot2^t$, $0\le t\le8$. Left: the simple upper floors $U_{\rm T}(q+1)$ and $U_{\rm e}(q+2)$ minus the new odd/even counts remain nine/ten; the corresponding gaps to $\G$ grow. The vertical axis is logarithmic. Right: new counts minus $\G$, equal to $20\cdot2^t-11$ and $10\cdot2^t-11$. The even gain is $-1$ at $t=0$, so the envelope retains 1220 at order 62. These upper comparisons concern simple arrangements.}\label{jfam:61-gaps}
\end{figure}

\begin{corollary}[Completed known orbits]\label{jfam:known}
The compatible seeds in Table~\ref{jfam:orbits} give unconditional odd/even Lean families. Both branches attain their classical simple upper floors.
\end{corollary}
\begin{table}[tb]
\centering\small
\begin{tabular}{@{}rrll@{}}
\toprule
Seed & $q$ & Count at $q+1$ & Count at $q+2$\\\midrule
$33:341$ & $32\cdot2^t$ & $(1024\,4^t-1)/3$ & odd count $+16\cdot2^t$\\
$37:431$ & $36\cdot2^t$ & $432\,4^t-1$ & odd count $+18\cdot2^t$\\
$49:767$ & $48\cdot2^t$ & $768\,4^t-1$ & odd count $+24\cdot2^t$\\
\bottomrule
\end{tabular}
\caption{Formal completion of known numerical orbits: Forge--Ram\'irez Alfons\'in, Parpalak--Utkin/Blanc and BBL, respectively~\cite{forge1998,parpalakutkin2026,blanc2011,bbl2008}.}\label{jfam:orbits}
\end{table}
The 33/49 finite seed checks were unfinished in the 3 October release; they are now discharged. The 37 orbit uses a compatible refit of Poola's $37:431$ input~\cite{poola2026} and the actual odd deficit-two successor, requiring no additional visibility check. The simple counts $38:449$ and $50:791$ are below retained nonsimple witnesses $38:450$ and $50:792$. The corrected five-line seed (five triangles, three caps, two visible pairs on $0<\varepsilon\le10^{-5}$) and the external uniform 19:107 seed check (17 caps on $0<\varepsilon\le10^{-3}$) retain their earlier evidence and attribution. The inherited odd families at $q=4\cdot2^t$ and $q=18\cdot2^t$ count $(q^2-1)/3$ and $q^2/3-1$, respectively; their simple even companions add $q/2$ by the now-verified deficit-two rule. The table completes the specified later tails in Lean.

Let $C(n)$ be the largest saved finite count, or zero. Retain
\begin{equation}\label{j:envelope}
 L(n)=\max\{\G(n),C(n)\}\le\K(n).
\end{equation}
Let $R(n)$ select the $q=10\cdot2^t$ odd/even counts, and vanish otherwise.
\begin{theorem}[Retained and enlarged envelopes]\label{j:new-envelope}
For every natural $n$,
\begin{equation}\label{j:new-envelope-formula}
 \K(n)\ge H(n):=\max\{L(n),R(n)\}\ge L(n).
\end{equation}
This is strict at $10\cdot2^{t+5}+1$ and $10\cdot2^{t+5}+2$. The enlarged executable envelope
\begin{equation}\label{j:portfolio}
 \K(n)\ge P(n):=\max\bigl(H(n),\max_{0\le t\le n}
 \{F_{33}(n,t),F_{37}(n,t),F_{49}(n,t),F_{61}(n,t)\}\bigr)
\end{equation}
retains $H$ pointwise, where candidates vanish away from their specified odd/even orders.
\end{theorem}
\begin{proof}
Take the finite maximum of actual witnesses. The old catalog stops at order 195. Beyond that,
\begin{equation}\label{j:new-gains}
 T_q-\G(q+1)=\floor{q/3}-2,\quad
 T_q+q/2-\G(q+2)=q/2-\floor{q/3}-2,
\end{equation}
both positive from $q=320$. The new maximum is
\proofref{Kobon/OpenMathConstructionEnvelope.lean}{97}{\lean{all_n}};
\proofref{Kobon/OpenMathConstructionEnvelope.lean}{100}{\lean{previous_le}}
proves retention. Its leading quadratic quality is inherited.
\end{proof}
The old checkpoints $81:2132$, $82:2172$, $161:8532$, $162:8612$, and existence bounds $321:34132$, $322:34292$, $641:136532$, $642:136852$ remain. Their gains at 321/322 are 104/52, at 641/642 are 211/105. Existing small-order certificates may dominate the simple families, including $12:38$, $21:133$, $22:143$, $41:533$ and $42:553$. No new theorem retroactively certifies unfinished archived coordinate files.
